% ============================================================================
\chapter{Entender los resultados}\label{cap:resultados}
% ============================================================================

Este capítulo explica, con los números de ejemplo del libro, qué significa cada indicador y qué decisiones ayuda a tomar.

\section{De los ingresos a la utilidad}

\begin{figure}[H]\centering
\begin{tikzpicture}[
  b/.style={draw=azul,minimum height=0.9cm,text width=6.3cm,align=left,font=\small,inner sep=4pt},
  r/.style={b,fill=azulclaro,font=\small\bfseries}]
\node[r] (i) {Ingresos = kg $\times$ tarifa de cada región};
\node[b,below=0pt of i] (v) {$-$ Costos variables (puerto, transitaria, distribución, comisión variable, reclamaciones, banca)};
\node[r,below=0pt of v] (m) {= Margen de contribución};
\node[b,below=0pt of m] (f) {$-$ Costos fijos del mes (o la parte de cada contenedor)};
\node[b,below=0pt of f] (t) {$-$ Tributos sobre ingresos (10\,\% + 1\,\%)};
\node[r,below=0pt of t] (u) {= Utilidad antes de impuesto (UAI)};
\node[b,below=0pt of u] (iu) {$-$ Impuesto sobre utilidades (35\,\%, estimado)};
\node[r,below=0pt of iu,fill=green!15] (n) {= Utilidad neta estimada};
\node[right=0.6cm of m,text width=6cm,font=\small,align=left] {Con los datos de ejemplo:\\ingresos \vIngtotal; margen de contribución \vMCpct{} de los ingresos.};
\node[right=0.6cm of u,text width=6cm,font=\small,align=left] {UAI \vUAItotal};
\node[right=0.6cm of n,text width=6cm,font=\small,align=left] {Utilidad neta \vUNtotal};
\end{tikzpicture}
\caption{Cómo se forma la utilidad}
\end{figure}

\section{Los indicadores, uno por uno}

\begin{description}[style=nextline,leftmargin=1.2em,itemsep=6pt]
\item[Margen de contribución (MC)]
  Lo que queda de los ingresos después de pagar los gastos que dependen de cada contenedor. Es el dinero disponible para pagar los gastos fijos y
  para ganar. Un MC del \vMCpct{} quiere decir que ese porcentaje de cada dólar que entra queda disponible para cubrir los gastos fijos y dar ganancia.
\item[Utilidad antes de impuesto (UAI)]
  La ganancia del período después de todos los gastos y de los tributos sobre los ingresos. Si es negativa (en rojo y entre paréntesis), el
  período tuvo pérdidas.
\item[Utilidad neta estimada]
  La UAI menos el impuesto sobre utilidades estimado (35\,\%). El impuesto real se liquida \textbf{una vez al año}: use la tabla anual de \hoja{RESUMEN\_MES}
  y confírmela con el contador.
\item[Ingreso por kg y costo total por kg]
  Cuánto entra y cuánto cuesta, en promedio, cada kg transportado. En el ejemplo entran \vIngKg{} y cuesta \vCostoKg{} por kg. La diferencia es la
  utilidad por kg. Si el costo por kg sube de un mes a otro, busque la causa en la estructura de costos del \hoja{DASHBOARD}.
\item[Punto de equilibrio]
  Los kg (o contenedores) necesarios para cubrir \textbf{exactamente} todos los gastos fijos: por debajo hay pérdidas; por encima, ganancias.
  En el ejemplo hacen falta \vPEkg{} kg en el período.
\item[Margen de seguridad]
  Cuánto pueden bajar los kg antes de entrar en pérdidas. Un \vMSeg{} significa que el volumen podría caer ese porcentaje y la empresa todavía
  no perdería. Cuanto mayor, más tranquilidad.
\item[Fijos no absorbidos]
  Gastos fijos que ningún contenedor cubrió: un mes sin arribos, o una región sin carga ese mes. Indican \textbf{capacidad ociosa}, es decir,
  se pagó una estructura que no se usó.
\item[Desvío del combustible]
  Cuánto se apartó el consumo real de la norma. Positivo: se gastó más de lo normal.
\end{description}

\section{Leer la rentabilidad por región}

La utilidad de cada región se calcula así: sus ingresos, menos sus propios costos (combustible, desagrupe regional, comisión y gastos fijos de la
región), menos \textbf{su parte de los costos compartidos} (puerto, transitaria y gastos fijos generales) según sus kg.

\begin{ejemplo}
Con los datos de ejemplo, Occidente gana \vUAIocc, Centro \vUAIcen{} y Oriente \textcolor{red}{\vUAIori}. Oriente cobra lo mismo por kg, pero tiene
dos tramos de transporte y un segundo desagrupe. Con una tarifa única, \textbf{Occidente está pagando las pérdidas de Oriente}.
El \hoja{SIMULADOR} calcula la tarifa mínima de Oriente (\vSimTminOri{} por kg) y el ejercicio~7 muestra el efecto de subirla.
\end{ejemplo}

\section{Preguntas que el libro le ayuda a responder}

\begin{center}
\begin{tabularx}{\linewidth}{@{}>{\raggedright\arraybackslash}p{6.2cm}>{\raggedright\arraybackslash}X@{}}
\toprule
\textbf{Pregunta} & \textbf{Dónde mirar} \\ \midrule
¿Ganamos dinero este mes? & \hoja{RESUMEN\_MES}, columna «Utilidad antes de impuesto» \\
¿Cuál fue el contenedor menos rentable y por qué? & \hoja{ENTRADA}, columna AC; detalle en \hoja{CALCULO} \\
¿Qué región conviene revisar? & \hoja{DASHBOARD}, rentabilidad por región \\
¿Qué tarifa mínima debemos cobrar? & \hoja{SIMULADOR}, fila «Tarifa mínima de equilibrio» \\
¿Cuántos contenedores necesitamos al mes? & \hoja{SIMULADOR}, punto de equilibrio y tabla de contenedores \\
¿Qué pasa si sube el diésel? & \hoja{SIMULADOR}, tabla de sensibilidad \\
¿En qué se va el dinero? & \hoja{DASHBOARD}, estructura de costos \\
¿Cuánto pagaremos de impuesto este año? & \hoja{RESUMEN\_MES}, liquidación anual estimada \\
\bottomrule
\end{tabularx}
\end{center}

\section{Cómo presentar los resultados a la dirección}
\begin{pasos}
\paso En \hoja{DASHBOARD}, elija el período.
\paso Exporte la hoja a PDF: \menu{Archivo $\rightarrow$ Exportar $\rightarrow$ Crear PDF}, o \menu{Imprimir} y elija la impresora «PDF».
\paso Acompáñelo con tres frases: la utilidad del período y su tendencia, la región más débil y su causa, y una propuesta (tarifa, consolidar viajes, controlar combustible\dots).
\end{pasos}
